\chapter{From Steps to the Archimedean Spiral}\label{sec:spirale}

%═══════════════════════════════════════════════════════════════════════════
\section{Geometric Motivation: From the Discrete Step Pattern to the
Continuous Spiral}
%═══════════════════════════════════════════════════════════════════════════

\subsection{Starting Point and Guiding Question}

In Chapter~\ref{sec:arithmetik} it was shown that the quadratic sequence
\[
  k(n) = \frac{2n^2 + 3n + 1}{6}
\]
takes integer values precisely for the index values $n \equiv 1, 5 \pmod{6}$.
The resulting discrete point set
\[
  \mathcal{P} = \bigl\{(n,\, k(n)) \mid n \equiv 1, 5 \pmod{6},\ n \geq 1\bigr\}
\]
has a parabolic shape in Cartesian coordinates, since $k(n)$ is a polynomial
of degree two in $n$. For the objective of this work --- the embedding of this
discrete structure into a geometrically regular, three-dimensional space --- the
Cartesian representation is, however, not suitable. It admits no natural
periodisation, no angular structure, and it does not make the underlying modular
regularity of the index $n$ geometrically visible.

\begin{tcolorbox}[
  title={Guiding question of this chapter},
  colback=gray!5,
  colframe=gray!50!black,
  enhanced,
  breakable
]
\emph{Does there exist a partition of the difference sequence of the
sequence $k(n)$ into equally-sized blocks such that the resulting block
sums can be interpreted as circumferences of concentric circles and the
resulting radius sequence defines an Archimedean spiral with a closed-form
slope constant? And if so: which block sizes admit it?}
\end{tcolorbox}

The answer, proved below: the block sums form an arithmetic progression
precisely for the even block sizes, and $\ell = 2$ is the minimal
admissible choice; combined with the circumference interpretation this
determines the spiral completely.

%─────────────────────────────────────────────────────────────────────────
\subsection{The Difference Sequence of the Integer Values}
\label{subsec:differenzstruktur}

We denote by $(n_k)_{k \geq 1}$ the ascending sequence of all
indices with $n_k \equiv 1, 5 \pmod{6}$:
\[
  n_1 = 1,\quad n_2 = 5,\quad n_3 = 7,\quad n_4 = 11,\quad
  n_5 = 13,\quad n_6 = 17,\quad n_7 = 19,\quad n_8 = 23,\quad \ldots
\]
The corresponding function values (integer values of the sequence) are:
\[
  p_k = k(n_k)\colon\quad 1,\ 11,\ 20,\ 46,\ 63,\ 105,\ 130,\ 188,\ 221,\ldots
\]
We define the \emph{first difference sequence} of this value sequence:
\[
  \delta_k = p_{k+1} - p_k = k(n_{k+1}) - k(n_k),\qquad k \geq 1.
\]
The first terms of this difference sequence are:
\[
  \delta_1 = 10,\quad \delta_2 = 9,\quad \delta_3 = 26,\quad
  \delta_4 = 17,\quad \delta_5 = 42,\quad \delta_6 = 25,\quad
  \delta_7 = 58,\quad \delta_8 = 33,\quad \ldots
\]

\begin{remark}[Alternating growth of the differences]
The difference sequence $(\delta_k)$ exhibits a clear alternating pattern:
\begin{itemize}
  \item The terms at \emph{odd} positions ($\delta_1, \delta_3, \delta_5, \ldots$)
        take the values $10, 26, 42, 58, \ldots$ and form an arithmetic subsequence
        with common difference $16$.
  \item The terms at \emph{even} positions ($\delta_2, \delta_4, \delta_6, \ldots$)
        take the values $9, 17, 25, 33, \ldots$ and form an arithmetic subsequence
        with common difference $8$.
\end{itemize}
This pattern directly reflects the alternation of the index differences
$n_{k+1} - n_k$, which alternate between $\Delta n = 4$ (for
$n_k \equiv 1 \pmod{6}$, \emph{major steps}) and $\Delta n = 2$ (for
$n_k \equiv 5 \pmod{6}$, \emph{minor steps}).
\end{remark}

%─────────────────────────────────────────────────────────────────────────
\subsection{The Block-Formation Problem and Its Solution}
\label{subsec:gruppenbildung}

We pose the following algebraic question: which partition of the
difference sequence $(\delta_k)$ into consecutive, equally-sized blocks
of length $\ell \in \{1, 2, 3, \ldots\}$ produces a sequence of block
sums that forms an \emph{arithmetic progression}?

Formally, for a block size $\ell \geq 1$ we define the sequence of
block sums
\[
  S_m^{(\ell)} = \sum_{j=1}^{\ell} \delta_{\ell(m-1)+j},\qquad m \geq 1,
\]
and ask: for which values of $\ell$ is $(S_m^{(\ell)})_{m \geq 1}$
an arithmetic sequence?

\begin{theorem}[Even block sizes and minimality of the pairing]\label{satz:eindeutigkeit}
The sequence of block sums $(S_m^{(\ell)})_{m \geq 1}$ is an arithmetic
progression if and only if $\ell$ is even. For $\ell = 2m$ the common
difference is
\[
  S_{i+1}^{(\ell)} - S_i^{(\ell)} = 24m^2 = 6\ell^2 .
\]
In particular, $\ell = 2$ is the minimal admissible block size, with
common difference $24$.
\end{theorem}

\begin{proof}
By Theorem~\ref{satz:minor_major_formeln} the difference sequence is
$\delta_{2t-1} = 16t - 6$ (major) and $\delta_{2t} = 8t + 1$ (minor).

\textbf{Even $\ell = 2m$:} shifting an index by $\ell$ preserves its
parity and advances the period counter $t$ by $m$, so
\[
  \delta_{j+\ell} - \delta_j =
  \begin{cases}
    16m, & j \text{ odd},\\
    8m,  & j \text{ even}.
  \end{cases}
\]
Each block of length $2m$ contains exactly $m$ odd and $m$ even indices;
hence
\[
  S_{i+1}^{(\ell)} - S_i^{(\ell)}
  = \sum_{j \in \text{Block}_i} \bigl(\delta_{j+\ell} - \delta_j\bigr)
  = m \cdot 16m + m \cdot 8m = 24m^2 = 6\ell^2,
\]
independent of $i$: the block sums form an arithmetic progression.

\textbf{Odd $\ell$:} now $j$ and $j+\ell$ have opposite parity. Writing
$\ell = 2m+1$, a direct computation from the two closed formulae gives
\[
  \delta_{j+\ell} - \delta_j =
  \begin{cases}
    -8t + 4\ell + 3, & j = 2t-1 \text{ odd},\\
    \phantom{-}8t + 8\ell + 1, & j = 2t \text{ even}.
  \end{cases}
\]
Since $\ell$ is odd, the parity of the first index of $\text{Block}_i$
alternates with $i$, so consecutive blocks contain the two term types in
the ratios $(m{+}1) : m$ and $m : (m{+}1)$ alternately. Summing over a
block yields
\[
  S_{i+1}^{(\ell)} - S_i^{(\ell)} = 6\ell^2 + (-1)^i\,(4\ell i + 3),
\]
an expression whose deviation from $6\ell^2$ alternates in sign and grows
without bound; $(S_m^{(\ell)})$ is therefore not an arithmetic
progression for any odd $\ell$. The minimality of $\ell = 2$ among the
even block sizes is immediate.
\end{proof}

\begin{remark}[Numerical illustration]
For $\ell = 2$: sums $19, 43, 67, 91$ with constant difference $24$.
For $\ell = 4$: sums $62, 158, 254, 350$ with constant difference
$96 = 6 \cdot 4^2$. For $\ell = 3$: sums $45, 84, 165$ with differences
$39, 81$ --- not arithmetic, in accordance with the closed form
$6\ell^2 + (-1)^i(4\ell i + 3) = 54 + (-1)^i(12i+3)$.
\end{remark}

\begin{remark}[Connection to the OEIS]
	The sequence of circumferences $U_m = 24m - 5$ is listed in the OEIS
	under A350051~\cite{OEIS} as part of the trisection of
	$A017101 = \{3 + 8k\}_{k \geq 0}$. It is congruent to
	$1 \pmod{6}$ --- in agreement with the congruence structure
	of the sequence $k(n)$.
\end{remark}

\begin{remark}[Compatibility of the minimal pairing]
Every even block size yields an arithmetic progression
(Theorem~\ref{satz:eindeutigkeit}); the pairing $\ell = 2$ is
distinguished twice over. It is minimal, and it is the only even block
size compatible with the angular assignment of the construction: one
block then comprises exactly one major and one minor step, i.e.\ the six
index units $\Delta n = 6$ of a full period, which Chapter~\ref{ch:helix}
maps to one full revolution $\Delta\theta = 2\pi$.
\end{remark}

%─────────────────────────────────────────────────────────────────────────
\subsection{Geometric Interpretation of the Block Sums as Circumferences}
\label{subsec:geometrische_interpretation}

Having established that the pairwise block sums $S_m^{(2)} = 24m - 5$ form
an arithmetic progression, we now present the decisive geometric argument.

An Archimedean spiral $r(\varphi) = a\varphi$ has the defining
property of \emph{constant radial difference}: for each complete
revolution $\Delta\varphi = 2\pi$ one has
\[
  \Delta r = 2\pi a = \text{const}.
\]
Equivalently: the circumferences $U_m = 2\pi R_m$ of successive
spiral windings themselves form an arithmetic progression with constant
difference $\Delta U = 2\pi \Delta R$.
\clearpage

\begin{tcolorbox}[
  title={Justification of the circumference identification},
  colback=blue!3,
  colframe=blue!40!black,
  enhanced,
  breakable
]
The integer points of the sequence $k(n)$ admit a polar arrangement in
which the two residue classes occupy fixed angular directions: points
with $n_k \equiv 1 \pmod{6}$ along one ray, points with
$n_k \equiv 5 \pmod{6}$ along another (made precise by the twin zones of
Chapter~\ref{ch:helix}). This suggests interpreting each pair of
consecutive steps --- one major and one minor --- as a complete
\emph{period} of the arrangement, corresponding to one winding.

Under this interpretation the sum of the increments within a period
corresponds to the circumference increment of a circle. The
identification itself is a definition
(Definition~\ref{def:kreisumfaenge}); Theorem~\ref{satz:eindeutigkeit}
then shows that the even pairings --- minimally $\ell = 2$ --- are
exactly those for which the sums grow arithmetically. This is the
algebraic statement that the identification is \emph{consistent}: with
the minimal pairing a spiral with \emph{constant} winding spacing
arises, and its consequences are forced, as the following sections show.
\end{tcolorbox}

\begin{definition}[Circumferences of the sequence]\label{def:kreisumfaenge}
We identify the pairwise block sums of the difference sequence
with the circumferences of a family of concentric circles:
\[
  U_m \coloneqq S_m^{(2)} = \delta_{2m-1} + \delta_{2m} = 24m - 5.
\]
\end{definition}

\begin{remark}[Classification by placement rule]\label{rem:placement_classification}
The circumference reading is one of several possible geometric
interpretations of the arithmetic data, and the choice classifies the
resulting picture --- not by curve type, but by the \emph{placement map}
$n \mapsto (r, \theta)$. Reading the linear block sums as circumferences
gives $r \propto m$: a constant occupation of six integer points per
winding, the construction of this work. Reading cumulative \emph{counts}
as areas gives $r \propto \sqrt{N}$: the placement of the Sacks
spiral~\cite{Sacks2003}, which carries all integers with $2n+1$ points
on the $n$-th winding --- its underlying curve is likewise Archimedean;
only the occupation density grows. A root radius combined with a
\emph{linear} angle ($r \propto \sqrt{n}$, $\theta \propto n$) yields
the Fermat spiral of Vogel's phyllotaxis model. Within the linear
branch, the circumference is singled out among the quantities
proportional to $R$ by a dimensional property: it is the one choice for
which the block sum has the same meaning as the quantity it is
identified with --- the length traversed per revolution, exactly for the
circle and to leading order for the spiral
(Chapter~\ref{sec:bogenlaenge_spirale}).
\end{remark}

%═══════════════════════════════════════════════════════════════════════════
\section{Derivation of the Major and Minor Step Formulae}
\label{sec:minor_major_herleitung}
%═══════════════════════════════════════════════════════════════════════════

\subsection{Explicit Computation of the Step Sizes}

We derive the closed-form formulae for the major and minor step sizes
directly from the definition of $k(n)$.

The closed-form step formulae were established arithmetically in
Chapter~\ref{sec:arithmetik} (Theorem~\ref{satz:minor_major_formeln},
Definition~\ref{def:minor_major}):
\[
  \Delta_{\mathrm{Major}}(m) = 16m - 6,\qquad
  \Delta_{\mathrm{Minor}}(m) = 8m + 1.
\]
The first periods:
\begin{table}[htbp]
\centering
\caption{Step sizes of the first periods}
\begin{tabular}{c|cc|cc|c}
\toprule
Period $m$ & Starting value $n$ & $\Delta_{\mathrm{Major}}$ &
Intermediate value $n$ & $\Delta_{\mathrm{Minor}}$ & $U_m$ \\
\midrule
1 & 1  & 10 & 5  & 9  & 19 \\
2 & 7  & 26 & 11 & 17 & 43 \\
3 & 13 & 42 & 17 & 25 & 67 \\
4 & 19 & 58 & 23 & 33 & 91 \\
5 & 25 & 74 & 29 & 41 & 115 \\
\bottomrule
\end{tabular}
\end{table}

\begin{remark}[Verification of the formulae]
For $m=1$: $\Delta_{\mathrm{Major}}(1) = 16\cdot1 - 6 = 10$,
$\Delta_{\mathrm{Minor}}(1) = 8\cdot1 + 1 = 9$.
For $m=2$: $\Delta_{\mathrm{Major}}(2) = 32-6 = 26$,
$\Delta_{\mathrm{Minor}}(2) = 17$. The formulae match the computed differences.
\end{remark}

\subsection{The Unified Generator Function}

The two subsequences can be combined into a single formula
that generates all steps of the interleaved sequence $(\delta_k)$.

\begin{theorem}[Closed-form generator of the difference sequence]
The complete difference sequence $(\delta_k)_{k \geq 1}$ in its natural
order starting from $n = 1$:
\[
  \delta_1 = 10,\; \delta_2 = 9,\; \delta_3 = 26,\; \delta_4 = 17,\;
  \delta_5 = 42,\; \delta_6 = 25,\; \delta_7 = 58,\; \delta_8 = 33,\; \ldots
\]
(\textbf{major, minor, major, minor, \ldots}) is generated by the closed-form
formula
\begin{equation}
  a_k = \frac{1}{2}\bigl[(12k+3) + (-1)^{k+1}(4k+1)\bigr],
  \qquad k \geq 1,
\end{equation}
where $a_1 = 10$ corresponds to the first major step (from $n_1 = 1$ to
$n_2 = 5$).
\end{theorem}

\begin{proof}
For $k = 2m-1$ (odd positions, major steps):
\begin{align*}
  a_{2m-1}
  &= \frac{1}{2}\bigl[(12(2m-1)+3) + (-1)^{2m}(4(2m-1)+1)\bigr] \\
  &= \frac{1}{2}\bigl[(24m-9) + (8m-3)\bigr]
  = \frac{32m-12}{2} = 16m-6.
\end{align*}
This is exactly $\Delta_{\mathrm{Major}}(m) = 16m-6$. \quad
Check: $a_1 = 16\cdot1-6 = 10$. \checkmark

For $k = 2m$ (even positions, minor steps):
\begin{align*}
  a_{2m}
  &= \frac{1}{2}\bigl[(12(2m)+3) + (-1)^{2m+1}(4(2m)+1)\bigr] \\
  &= \frac{1}{2}\bigl[(24m+3) - (8m+1)\bigr]
  = \frac{16m+2}{2} = 8m+1.
\end{align*}
This is exactly $\Delta_{\mathrm{Minor}}(m) = 8m+1$. \quad
Check: $a_2 = 8\cdot1+1 = 9$. \checkmark

The period sums remain identical:
\[
  U_m = a_{2m-1} + a_{2m}
      = (16m-6) + (8m+1)
      = 24m - 5. \qquad\checkmark
\]
\end{proof}

%═══════════════════════════════════════════════════════════════════════════
\section{From Steps to Circumferences}
%═══════════════════════════════════════════════════════════════════════════

The step notation $\Delta_{\mathrm{Minor},m}$, $\Delta_{\mathrm{Major},m}$ was introduced in Definition~\ref{def:minor_major}.

By Definition~\ref{def:kreisumfaenge}, one major and one minor step
together form the $m$-th circumference; by
Theorem~\ref{satz:eindeutigkeit} (case $\ell = 2$),
\begin{equation}
  U_m = \Delta_{\mathrm{Major},m} + \Delta_{\mathrm{Minor},m}
      = (16m-6) + (8m+1) = 24m - 5,
  \qquad U_{m+1} - U_m = 24.
\end{equation}

%═══════════════════════════════════════════════════════════════════════════
\section{Radius Sequence}
%═══════════════════════════════════════════════════════════════════════════

From the circumferences the radius sequence follows via $U_m = 2\pi R_m$:
\begin{equation}
  R_m = \frac{U_m}{2\pi} = \frac{24m - 5}{2\pi}.
\end{equation}

\begin{theorem}[Constant radius difference --- key result]
\begin{equation}
  \boxed{\Delta R = R_{m+1} - R_m = \frac{24}{2\pi} = \frac{12}{\pi}
  \approx 3.819718634205}
\end{equation}
\end{theorem}

\begin{proof}
Direct computation:
\[
  R_{m+1} - R_m
  = \frac{24(m+1)-5}{2\pi} - \frac{24m-5}{2\pi}
  = \frac{24}{2\pi}
  = \frac{12}{\pi}.
\]
\end{proof}

%═══════════════════════════════════════════════════════════════════════════
\section{Archimedean Spiral}
%═══════════════════════════════════════════════════════════════════════════

\begin{definition}
An \emph{Archimedean spiral} is a plane curve in polar coordinates
of the form
\[
  r(\theta) = a\theta,
\]
with slope constant $a > 0$. The characterising property of this
class of curves is the constant radial difference between successive
windings: for each complete revolution $\Delta\theta = 2\pi$,
\[
  \Delta r = r(\theta + 2\pi) - r(\theta) = 2\pi a.
\]
\end{definition}

The minimality of the pairing proved in Theorem~\ref{satz:eindeutigkeit}
and the resulting constant $\Delta R = 12/\pi$ determine
the slope constant $a$:
\begin{align}
  2\pi a &= \Delta R = \frac{12}{\pi},\\[4pt]
  a &= \frac{12}{2\pi^2} = \frac{6}{\pi^2}.
\end{align}

\begin{theorem}[Spiral equation]
\begin{equation}
  \boxed{
    r(\theta) = \frac{6}{\pi^2}\,\theta
    \approx 0.607927101854 \cdot \theta
  }
\end{equation}
\end{theorem}

\begin{remark}[Connection to the Basel problem]
The slope constant
\[
  a = \frac{6}{\pi^2} = \frac{1}{\zeta(2)},
  \qquad
  \zeta(2) = \sum_{n=1}^{\infty} \frac{1}{n^2} = \frac{\pi^2}{6}
\]
is the reciprocal of the Riemann zeta function at $s=2$, i.e.\ the
solution of the classical Basel problem. The value $\zeta(2) = \pi^2/6$
was proved by Leonhard Euler (published 1740)~\cite{Euler1740,Hardy2008}. The fact
that the spiral constant, obtained from purely combinatorial-arithmetic
considerations on the difference structure of the sequence $k(n)$,
coincides with $1/\zeta(2)$ is a structural identity.
\end{remark}

\begin{remark}[The spiral constant as the mean coprimality density]
\label{rem:euler_coprimality}
The coincidence of the spiral constant $a = 6/\pi^2$ with the reciprocal
of $\zeta(2)$ carries a deeper number-theoretic meaning that extends well
beyond the Basel problem.

A classical result of analytic number theory states that the average
behaviour of Euler's $\varphi$-function satisfies
\begin{equation}\label{eq:euler_mean}
  \frac{1}{x}\sum_{n \leqslant x} \frac{\varphi(n)}{n}
  \;\longrightarrow\; \frac{6}{\pi^2}
  \qquad (x \to \infty)
\end{equation}
(see, e.g.,~\cite{Hardy2008}).
The ratio $\varphi(n)/n = \prod_{p \mid n}(1 - 1/p)$ measures the
\emph{proportion of integers up to~$n$ that are coprime to~$n$}.
The limit $6/\pi^2$ is therefore the \emph{mean coprimality density}
of the natural numbers: if one draws a random pair $(a,n)$ with
$1 \leqslant a \leqslant n$, the probability that $\gcd(a,n) = 1$
is asymptotically equal to $6/\pi^2$.

The spiral constant of the present work thus admits three simultaneous
interpretations:
\begin{enumerate}
\item $1/\zeta(2)$ — the reciprocal of the series $\sum 1/n^2$
      \emph{(Basel problem, Euler 1734)};
\item the mean coprimality density $\overline{\varphi(n)/n}$
      \emph{(summatory number theory)};
\item the slope~$a$ of the Archimedean spiral, which follows
      \emph{uniquely} from the difference structure of~$k(n)$
      \emph{(present work)}.
\end{enumerate}

This threefold connection is not coincidental: the sequence $k(n)$
takes integer values precisely when $n \equiv 1,5\pmod{6}$, that is,
for those~$n$ that are coprime to~$6$. The proportion of these
congruence classes is $\varphi(6)/6 = 2/6 = 1/3$, and the spiral
constant $a = 6/\pi^2$ encodes, in a sense, the \emph{universal
extension} of this local coprimality ratio to all primes simultaneously.
\end{remark}

\begin{remark}[Distribution function of the error term and the
Erd\H{o}s--Shapiro theorem]
\label{rem:erdos_shapiro}
Erd\H{o}s and Shapiro~\cite{ErdosShapiro1955} studied the error function
\begin{equation}\label{eq:H_erdos}
  H(x) = \sum_{n \leqslant x} \frac{\varphi(n)}{n}
         - \frac{6}{\pi^2}\,x,
\end{equation}
which measures the deviation of the cumulative values of $\varphi(n)/n$
from the linear main term $6x/\pi^2$. Their principal result is that
$H(x)$ possesses a \emph{continuous distribution function}: writing
$N(n,u) = \#\{m \leqslant n : H(m) \geqslant u\}$, the limit
\[
  F(u) = \lim_{n \to \infty} \frac{N(n,u)}{n}
\]
exists for every~$u$, and $F$ is continuous.

This result is relevant to the present work for two reasons:
\begin{enumerate}
\item The constant $6/\pi^2$ about which~$H(x)$ oscillates is
      \emph{exactly} the spiral constant~$a$. The Erd\H{o}s--Shapiro
      theorem shows that the fluctuations around this main term are
      statistically regular — the coprimality density is thus robust
      not only on average, but also in a probabilistic sense.
\item The existence and continuity of the distribution function of~$H(x)$
      relies crucially on the continuity of the distribution function
      of $\varphi(n)/n$ itself (Erd\H{o}s--Wintner~\cite{ErdosWintner1939}).
      The primes $p > 3$ satisfy $\varphi(p)/p = 1 - 1/p \to 1$ and
      therefore cluster at the upper end of this distribution. On the
      conical helix they appear as a geometrically distinguished
      subset — the statistical distribution of $\varphi(n)/n$ thereby
      acquires a spatial interpretation.
\end{enumerate}

The present construction thus \emph{geometrises} an arithmetic
quantity whose statistical fine structure has been the subject of
deep analytic investigations~\cite{ErdosShapiro1955,
Erdos1974distribution, ErdosHall1973}.
\end{remark}

\begin{remark}[Visible lattice points and squarefree numbers]
\label{rem:visible_squarefree}
Two further classical interpretations of the constant $6/\pi^2$ deserve
mention, since both are geometric in nature and thus close to the spirit
of the present construction (see~\cite{Hardy2008}).

First, a lattice point $(x,y) \in \mathbb{Z}^2$ is called \emph{visible}
from the origin if no other lattice point lies on the open segment
between the origin and $(x,y)$; this holds precisely when
$\gcd(x,y)=1$. The density of visible points among all lattice points is
\begin{equation}
  \lim_{R \to \infty}
  \frac{\#\{(x,y) : \|(x,y)\| \leqslant R,\ \gcd(x,y)=1\}}
       {\#\{(x,y) : \|(x,y)\| \leqslant R\}}
  = \frac{6}{\pi^2}.
\end{equation}

Second, the density of the squarefree integers is likewise
\begin{equation}
  \lim_{x \to \infty} \frac{\#\{n \leqslant x : n \text{ squarefree}\}}{x}
  = \prod_{p} \Bigl(1 - \frac{1}{p^2}\Bigr)
  = \frac{1}{\zeta(2)} = \frac{6}{\pi^2},
\end{equation}
where the Euler product on the left is term by term the same product
that yields the spiral constant~$a$ in the present work.

The list of interpretations of the slope~$a$ thus extends to five: the
reciprocal of $\zeta(2)$, the mean coprimality density
(Remark~\ref{rem:euler_coprimality}), the density of visible lattice
points, the density of squarefree integers, and the slope of the
Archimedean spiral derived here. The visible-point interpretation is
noteworthy in the present context because it is itself a statement about
straight lines through integer points of the plane --- the same
raw material from which the spiral is built. Whether this analogy can be
made structural is an open question of this work, not an established
result.
\end{remark}

\begin{remark}[Farey fractions and the Franel--Landau criterion]
\label{rem:farey_franel}
The Farey sequence $F_N$ consists of all reduced fractions in $[0,1]$
with denominator at most~$N$. Its cardinality satisfies
\begin{equation}\label{eq:farey_count}
  |F_N| = 1 + \sum_{n \leqslant N} \varphi(n)
        \;\sim\; \frac{3}{\pi^2}\,N^2
        \;=\; \frac{a}{2}\,N^2 ,
\end{equation}
with the same coefficient $a/2$ that appears in the arc-length
asymptotics $s(\theta) \approx \tfrac{a}{2}\theta^2$ of the spiral
(Chapter~\ref{sec:bogenlaenge_spirale}). Both statements are counting
results governed by the mean of $\varphi$; the identical quadratic
form is recorded here as a formal correspondence, without a claim
that the two counts measure the same quantity.

The connection to the Riemann zeta function is provided by the
Franel--Landau criterion~\cite{Franel1924, Landau1924}: writing
$f_1 < f_2 < \cdots < f_{|F_N|}$ for the elements of $F_N$ and
$\delta_j = f_j - j/|F_N|$ for their deviations from perfect
equidistribution, the Riemann hypothesis is \emph{equivalent} to
\begin{equation}
  \sum_{j=1}^{|F_N|} |\delta_j| = O\bigl(N^{1/2+\varepsilon}\bigr)
  \qquad \text{for every } \varepsilon > 0.
\end{equation}
The quality of the equidistribution of the reduced fractions --- the
same arithmetic objects whose count carries the coefficient $a/2$ ---
therefore encodes the position of the nontrivial zeros of
$\zeta(s)$.
\end{remark}

\begin{remark}[Mertens' theorem and the Euler--Mascheroni constant]
\label{rem:mertens_gamma}
The spiral constant is the convergent second-power Euler product
$a = \prod_p (1 - p^{-2})$. Its first-power counterpart does not
converge to a positive limit; instead, Mertens'
theorem~\cite{Mertens1874} states
\begin{equation}\label{eq:mertens}
  \prod_{p \leqslant x} \Bigl(1 - \frac{1}{p}\Bigr)
  \;\sim\; \frac{e^{-\gamma}}{\ln x}
  \qquad (x \to \infty),
\end{equation}
where $\gamma = 0.5772156649\ldots$ is the Euler--Mascheroni constant.
The partial product on the left is the density of integers free of
prime factors $\leqslant x$; for $x = 3$ it equals
$(1-\tfrac12)(1-\tfrac13) = \tfrac13 = \varphi(6)/6$, the local
coprimality ratio of the congruence classes $n \equiv 1,5 \pmod 6$ on
which the sequence $k(n)$ is integer-valued.

The two products thus continue the same local datum in two different
ways: with squared factors the product converges and yields the spiral
constant~$a$; with first powers it decays logarithmically, and
$\gamma$ is the constant that calibrates this decay. The constant
$\gamma$ is tied to $\zeta$ at its pole through the Laurent expansion
$\zeta(s) = \frac{1}{s-1} + \gamma + O(s-1)$, while $a$ evaluates
$1/\zeta$ at $s = 2$; the pair $(\gamma, a)$ therefore probes the same
function at its two arithmetically distinguished points.
\end{remark}


